\documentclass{article}

\begin{document}
    \begin{center}
        \Large \textbf{Rabbit Challenge 3: Backups, Migration, Testysteme} \\ [6pt]
        \normalsize
        Autor: Maximilian Jakob Maag B.Sc. \\
        Version: 1.0
    \end{center}

    \section{Intro}
    Backups und Desaster Recovery sind im Notfall unverzichtbar, werden aber eher nachlässig angelegt.
    In dieser Challenge geht es darum, ein Backup einer Nextcloud Installation zu erstellen und dieses Backup auf einem anderen Server wiederherzustellen.

    \section{Aufgaben}
    Ziel der Übung ist das Erstellen eines Backups und das Wiederhestellen auf einer anderen VM.

    \subsection{Aufgabe 1}
    Deployen Sie eine Instanz Ihres Nextcloud Playbooks in der Region US-eeast1.

    \subsection{Aufgabe 2}
    Die Geschäftsleitung teilt Ihnen mit, dass die Datenbank Ihrer Anwendung regelmäßig Entwicklern zu Testzwecken bereit gestellt werden muss.
    Schreiben Sie ein Playbook, dass ein Backup der Datenbank erstellt und diesen in ein S3 Bucket ablegt, wo das Entwickler Team es verwenden kann.

    \subsection{Aufgabe 3}
    Führen Sie ein vollständiges Backup der Nextcloud durch.
    Das blose Klonen der VM ist zu teuer. Sie müssen die Nextcloud, Alle Daten der Applikation und deren Einstellungen speichern ohne die VM zu klonen.
    Entwickeln Sie ein Playbook für diesen Use Case.

    \subsection{Aufgabe 4}
    Die Nextcloud muss aufgrund von GDPR Vorgaben aus der Region US-east1 in die Region EU-west1 migriert werden.
    Erstellen Sie ein Playbook, dass die Nextcloud in der Region EU-west1 wiederherstellt.
    Darüber hinaus muss die Nextcloud unter der selben URL erreichbar bleiben.
    Daten und Einstellungen dürfen nicht verloren gehen.

    \subsection{Aufgabe 5}
    Die Anforderungen an die Nextcloud sind gestiegen. Die Nutzerbasis hat sich durch den Zukauf einer neuen Tochtergesellschaft verdoppelt.
    Eine einzelne VM ist den Anforderungen nicht mehr gewachsen.
    Zwei mögliche Lösungen liegen auf dem Tisch:
    \begin{itemize}
        \item Die Nextcloud wird auf zwei VMs repliziert. Ein Load Balancer verteilt die Anfragen auf beide VMs.
        \item Die Nextcloud wird in einem Kubernetes Cluster betrieben, um die Skalierung zu erleichtern.
    \end{itemize}
    Setzen Sie eines der beiden Szenarien um.

    \subsection{Aufgabe 6}
    Das Entwickler Team möchte regelmäßig Testinstanzen der Nextcloud bereitgestellt bekommen.
    Erstellen Sie ein Playbook, dass eine Testinstanz der Nextcloud bereitstellt.
    Das blose Klonen der VM ist nicht ausreichend. Die Testinstanz muss unter der der URL test.nextcloud.example.com erreichbar sein und 
    die Datenbank muss mit einem Backup der Produktionsinstanz gefüllt werden.
    Außerdem müssen alle Benutzer und Einstellungen der Produktionsinstanz in der Testinstanz vorhanden sein.
    Die Testinstanz soll bei der Ausführung einer CI Pipeline deployed werden.
    Innerhalb der Pipeline soll nach dem deployen der Testinstanz eine suit mit E2E tests sicherstellen,
    dass Änderungen am firmen eigenen Nextcloud Plugin die Funktionalität der Nextcloud nicht beeinträchtigen.
    Nach dem Ausführen der Tests soll die Testinstanz automatisch wieder gelöscht werden.

\end{document}